\documentclass{amsart}
\usepackage{amsmath,amssymb,amsthm,hyperref}
\usepackage{algorithm}
\usepackage{algpseudocode}

\theoremstyle{plain}
\newtheorem{theorem}{Theorem}
\newtheorem{lemma}{Lemma}
\newtheorem{proposition}{Proposition}
\newtheorem{corollary}{Corollary}
\theoremstyle{definition}
\newtheorem{definition}{Definition}
\newtheorem{remark}{Remark}

\newcommand{\PP}{\mathbb{P}}
\newcommand{\AA}{\mathbb{A}}
\newcommand{\cO}{\mathcal{O}}
\DeclareMathOperator{\Spec}{Spec}
\DeclareMathOperator{\Proj}{Proj}
\DeclareMathOperator{\trdeg}{trdeg}
\DeclareMathOperator{\Pic}{Pic}
\DeclareMathOperator{\dimkv}{\dim_k}

\begin{document}

\title{An algorithm for deciding parametrizability and computing a rational
parametrization of an algebraic surface}
\maketitle

\begin{abstract}
Let $k$ be an algebraically closed field of characteristic zero, given through
effective field arithmetic, and let $X\subset\PP^n_k$ be an irreducible
projective surface (an affine surface is handled by taking projective closure).
We give a complete algorithm that, from defining equations of $X$, decides
whether $k(X)$ is a purely transcendental extension of $k$ of transcendence
degree two (i.e.\ whether $X$ is \emph{rational}) and, in the affirmative case,
outputs an explicit rational parametrization $\varphi\colon\AA^2_k\dashrightarrow
X$, $\varphi\in k(s,t)^{n+1}$ (resp.\ $k(s,t)^n$ in the affine case).  The
algorithm is built from four ingredients: (i) resolution of singularities to a
smooth projective model; (ii) the Castelnuovo rationality criterion, made
effective through cohomology of pluricanonical sheaves computable by Gr\"obner
bases; (iii) the minimal-model program for surfaces, which constructively
reduces a rational surface either to $\PP^2$ or to a conic bundle over $\PP^1$;
and (iv) rational point search on conics with rational function-field
coefficients. We prove correctness and termination and give a complexity bound
that is elementary recursive in the input size, with the dominant cost being
that of the Gr\"obner-basis computations.
\end{abstract}

\section{Formalization and conventions}

Throughout, $k$ is an algebraically closed field of characteristic $0$.  We
assume $k$ is presented \emph{effectively}: its elements appearing in inputs and
intermediate data lie in a finitely generated subfield
$k_0=\mathbb{Q}(\alpha_1,\dots,\alpha_r)$ in which we can perform the four field
operations, decide equality, and factor univariate polynomials (the latter is
possible: $k_0$ is a finitely generated extension of $\mathbb{Q}$, hence its
univariate polynomial factorization is reducible to factorization over
$\mathbb{Q}$ together with primitive-element manipulations, all algorithmic).
Because $k$ is algebraically closed, whenever the algorithm must adjoin a root
of a univariate polynomial $f\in k_0[u]$ it does so by passing to
$k_1=k_0[u]/(g)$ for an irreducible factor $g\mid f$; this $k_1\subset k$ is
again effective.  We use the phrase ``compute in $k$'' to mean ``perform the
operation in the current effective subfield, enlarging it by such algebraic
adjunctions when a root is required.''  All such adjunctions occurring in the
algorithm are of bounded degree (made explicit below), so the cumulative field
extension stays effective.

\begin{definition}[Input size]\label{def:size}
The input is a finite set of homogeneous generators $f_1,\dots,f_m\in
k[x_0,\dots,x_n]$ of the saturated homogeneous ideal $I(X)$ of an irreducible
projective surface $X\subset\PP^n_k$.  Let $d=\max_i\deg f_i$ and let $H$ bound
the bit size (height) of the coefficients of the $f_i$, expressed in the
effective field $k_0$.  The \emph{input size} is $N=(n,m,d,H)$; complexity is
measured in $n,m,d,H$.  An affine surface $X^{\mathrm{aff}}\subset\AA^n_k$ with
ideal $J\subset k[y_1,\dots,y_n]$ is replaced by the projective closure
$X=\overline{X^{\mathrm{aff}}}\subset\PP^n_k$ (the saturation of the
homogenization of $J$ with respect to the irrelevant ideal), which is computable
by a Gr\"obner basis; a parametrization of $X$ restricts to one of
$X^{\mathrm{aff}}$ by setting $x_0=1$ and clearing denominators.  Hence we treat
only the projective case.
\end{definition}

\begin{definition}[Rationality]\label{def:rat}
$X$ is \emph{rational} (= parametrizable in the sense of the problem) if
$k(X)\cong k(s,t)$ as $k$-algebras, equivalently if there is a dominant rational
map $\PP^2_k\dashrightarrow X$.  Since $\dim X=2$ and $X$ is irreducible,
$\trdeg_k k(X)=2$ automatically; the content of rationality is that the
extension is \emph{purely} transcendental.
\end{definition}

We recall the two facts from the theory of algebraic surfaces on which the
algorithm rests.

\begin{theorem}[Castelnuovo's rationality criterion; char $0$]
\label{thm:castelnuovo}
Let $S$ be a smooth projective surface over $k$.  Then $S$ is rational if and
only if
\[
q(S):=h^1(S,\cO_S)=0
\qquad\text{and}\qquad
P_2(S):=h^0(S,\cO_S(2K_S))=0 ,
\]
where $K_S$ is a canonical divisor.  Both $q(S)$ and $P_2(S)$ are birational
invariants of smooth projective surfaces.
\end{theorem}

\begin{proof}
This is Castelnuovo's theorem; see Beauville, \emph{Complex Algebraic
Surfaces}, Thm.~V.1, or Barth--Hulek--Peters--Van~de~Ven, \emph{Compact Complex
Surfaces}, Ch.~VI.  The birational invariance of $q$ and of the plurigenera
$P_m=h^0(mK_S)$ for smooth projective surfaces is standard (loc.\ cit.,
Ch.~III); it is exactly what allows us to read these numbers off any smooth
model.  Castelnuovo's criterion holds over any algebraically closed field of
characteristic $0$ because the statement is geometric and the surface and its
relevant cohomology are defined over a finitely generated subfield of $k$ which
embeds into $\mathbb{C}$; rationality descends and ascends along the extension
$k_0\hookrightarrow\mathbb{C}$ since it is detected by the geometric invariants
$q,P_2$ which are insensitive to algebraically closed base change in
characteristic $0$ (flat base change for coherent cohomology).
\end{proof}

\begin{theorem}[Structure of rational surfaces via the MMP]
\label{thm:mmp}
Let $S$ be a smooth projective rational surface over $k$.  Running the
minimal-model program (repeatedly contracting $(-1)$-curves) terminates in a
smooth projective \emph{minimal} model $S_{\min}$ that is either
\begin{enumerate}
\item $\PP^2_k$, or
\item a Hirzebruch surface $\mathbb{F}_e=\PP(\cO_{\PP^1}\oplus\cO_{\PP^1}(e))$
for some $e\ge0$, $e\ne1$ (a $\PP^1$-bundle over $\PP^1$).
\end{enumerate}
Each contraction of a $(-1)$-curve is a birational morphism with an explicit
inverse (a blow-up), realised by explicit rational maps.  In either terminal
case $S_{\min}$ carries an explicit birational map to $\PP^2_k$.
\end{theorem}

\begin{proof}
The classification of minimal rational surfaces is classical
(Beauville, Ch.~IV; the statement is the Enriques--Castelnuovo classification of
relatively minimal rational surfaces).  A $(-1)$-curve is a smooth rational
curve $E$ with $E^2=-1$; by Castelnuovo's contractibility criterion it is
contractible to a smooth point, and the contraction is the inverse of a blow-up.
For a Hirzebruch surface $\mathbb{F}_e$ the bundle projection
$\mathbb{F}_e\to\PP^1$ has a section, and a fibre together with the section give
a birational map $\mathbb{F}_e\dashrightarrow\PP^1\times\PP^1\dashrightarrow\PP^2$
(explicit: $\mathbb{F}_e$ is obtained from $\PP^1\times\PP^1$ or from $\PP^2$ by
a sequence of explicit elementary transformations / blow-ups, all invertible by
explicit rational maps).  $\PP^2$ is itself trivially birational to $\AA^2$.
\end{proof}

\begin{remark}
We do not need the full force of the Hirzebruch classification.  For the
algorithm it suffices to know that a smooth projective rational surface admits
an explicit birational map to a \emph{conic bundle} $C\to\PP^1$ with $C$ a
geometrically rational surface and generic fibre a conic, or directly to
$\PP^2$.  We use the adjunction/anticanonical machinery below to reach this form
constructively.
\end{remark}

\section{The algorithm}

\subsection{Overview}
The algorithm has the following stages.
\begin{enumerate}
\item[(A)] \textbf{Resolution.}  Compute a smooth projective model $S$ of $X$
together with an explicit birational map $\pi\colon S\dashrightarrow X$ and its
inverse.
\item[(B)] \textbf{Decision.}  Compute $q(S)$ and $P_2(S)$ via sheaf cohomology.
By Theorem~\ref{thm:castelnuovo}, $X$ is rational iff $q(S)=P_2(S)=0$.  If not,
output \textsc{Not parametrizable} and halt.
\item[(C)] \textbf{Reduction to a minimal model.}  When $S$ is rational, run the
constructive MMP (Theorem~\ref{thm:mmp}) to obtain $S_{\min}\in\{\PP^2,
\mathbb{F}_e\}$ and an explicit birational map $\beta\colon
S\dashrightarrow S_{\min}$.
\item[(D)] \textbf{Parametrization of the minimal model.}  Produce an explicit
birational map $\gamma\colon\PP^2\dashrightarrow S_{\min}$ (for the
$\mathbb{F}_e$ case via a conic bundle and a rational point search; for $\PP^2$
the identity).
\item[(E)] \textbf{Composition and dehomogenization.}  Set
$\varphi=\pi\circ\beta^{-1}\circ\gamma\colon\PP^2\dashrightarrow X$, compose the
explicit rational maps, restrict to the affine chart $\{[s:t:1]\}$ of the source,
and output $\varphi(s,t)$.
\end{enumerate}

We now make each stage effective and prove correctness.

\subsection{Stage A: resolution of singularities}

\begin{lemma}[Effective resolution]\label{lem:resolution}
There is an algorithm that, given $I(X)$ as in Definition~\ref{def:size},
outputs a smooth projective surface $S\subset\PP^{M}_k$ (for some computed $M$)
given by its homogeneous ideal, together with an explicit birational morphism
$\rho\colon S\to X$ given by homogeneous polynomials, an explicit rational
inverse $\rho^{-1}\colon X\dashrightarrow S$, and the smooth locus certificate.
\end{lemma}

\begin{proof}
Hironaka's resolution of singularities in characteristic $0$ is constructive and
has been turned into explicit algorithms (Bierstone--Milman; Villamayor;
Bravo--Encinas--Villamayor); for surfaces one may also use the classical
algorithm by repeated normalized blow-ups of the singular locus, which
terminates by Zariski's resolution for surfaces.  Concretely:
\begin{enumerate}
\item Compute the singular locus $\Sigma=\mathrm{Sing}(X)$ as the (radical of
the) ideal generated by $I(X)$ and the $3\times3$ minors of the Jacobian of a
generating set of $I(X)$ in each affine chart; this is a Gr\"obner-basis
computation.  $\Sigma$ has dimension $\le1$.
\item Blow up $X$ along $\Sigma$: form the Rees algebra
$\bigoplus_{j\ge0}\Sigma^j t^j$, present it by generators and relations
(elimination/Gr\"obner), and take $\Proj$.  This yields $X_1$ with a birational
morphism $X_1\to X$ given by explicit homogeneous maps; compute it as the closure
of the graph.
\item Normalize $X_1$ (compute the integral closure of its coordinate ring by
the algorithm of de Jong / Greuel--Pfister, again Gr\"obner based), obtaining
$X_1^{\nu}\to X_1$.
\item Repeat (1)--(3) on $X_1^{\nu}$.
\end{enumerate}
By Zariski's theorem on resolution of surface singularities (valid in char $0$),
the embedded multiplicity / Hilbert--Samuel data strictly decreases under the
normalized blow-up of the (reduced) singular locus, so after finitely many
steps $\mathrm{Sing}$ becomes empty; smoothness is certified by checking that the
Jacobian has full rank $3$ at every point, i.e.\ that the Jacobian ideal together
with $I$ is the irrelevant ideal (a Gr\"obner test for the empty vanishing
locus).  Each step is given by explicit homogeneous maps, so the composite
$\rho\colon S\to X$ and its rational inverse $\rho^{-1}$ are explicit.
\end{proof}

\subsection{Stage B: the decision (cohomology of $\cO_S$ and $2K_S$)}

\begin{lemma}[Effective cohomology and plurigenera]\label{lem:cohomology}
For a smooth projective $S\subset\PP^M_k$ given by its saturated homogeneous
ideal, the numbers $q(S)=h^1(\cO_S)$ and $P_2(S)=h^0(\cO_S(2K_S))$ are
computable.
\end{lemma}

\begin{proof}
\emph{Computing $h^1(\cO_S)$.}  Sheaf cohomology of a coherent sheaf on $\PP^M$
presented by a graded module is algorithmic: a finite free resolution of the
graded module $\bigoplus_d H^0(\PP^M,\cO_S(d))$ (equivalently of the homogeneous
coordinate ring $R=k[x]/I(X_{\mathrm{model}})$ as a module) is computed by
iterated syzygy (Gr\"obner) computations, and local cohomology / sheaf
cohomology is read off via the Bernstein--Gel'fand--Gel'fand correspondence or by
Eisenbud--Fl\o ystad--Schreyer's ``Tate resolution,'' all implemented as exact
linear-algebra over $k$ on the graded pieces.  Equivalently, $h^1(\cO_S)$ is
obtained from the Hilbert polynomial and the dimensions of finitely many graded
pieces of local cohomology modules $H^i_{\mathfrak m}(R)$, which are computable
from a free resolution of $R$.  Thus $q(S)$ is computable.

\emph{Computing the canonical sheaf.}  For a smooth surface $S\subset\PP^M$ with
$\dim S=2$, the canonical sheaf is
$\omega_S=\mathcal{E}xt^{M-2}_{\cO_{\PP^M}}(\cO_S,\omega_{\PP^M})$, with
$\omega_{\PP^M}=\cO_{\PP^M}(-M-1)$.  This $\mathcal{E}xt$ is computed from a free
resolution of $\cO_S$ over $\cO_{\PP^M}$ by dualizing and taking cohomology of the
dual complex; the result is a graded module $\Omega$ representing $\omega_S$,
output as a presentation.  All steps are Gr\"obner computations over $k$.

\emph{Computing $P_2$.}  The bicanonical sheaf $\cO_S(2K_S)=\omega_S^{\otimes2}$
is represented by the symmetric/tensor square module $\Omega^{(2)}$ (compute the
image of $\Omega\otimes\Omega\to(\omega_S^{\otimes2})$, i.e.\ a presentation of
the double-dual of $\Omega\otimes_R\Omega$; reflexive hull computed by a further
$\mathcal{H}om$ into $R$).  Then $P_2(S)=h^0(S,\omega_S^{\otimes2})=
\dim_k(\Omega^{(2)})_{\mathrm{deg}\,0}$ in the sheafified sense, again read off
from the BGG/Tate machine, or as
$h^0(S,\omega_S^{\otimes2})$ via the same local-cohomology computation applied
to $\Omega^{(2)}$.  Hence $P_2(S)$ is computable.
\end{proof}

\begin{proposition}[Correctness of the decision]\label{prop:decision}
$X$ is parametrizable (Definition~\ref{def:rat}) if and only if the smooth model
$S$ of Lemma~\ref{lem:resolution} satisfies $q(S)=0$ and $P_2(S)=0$.
\end{proposition}

\begin{proof}
$X$ is rational iff its smooth model $S$ is rational (rationality is a
birational invariant and $\rho\colon S\to X$ is birational).  By
Theorem~\ref{thm:castelnuovo}, $S$ is rational iff $q(S)=P_2(S)=0$.  Both numbers
are computed in Lemma~\ref{lem:cohomology}.  If either is nonzero the algorithm
correctly returns \textsc{Not parametrizable}.  (Note: the hypothesis of the
problem guarantees we are in the rational case, but the algorithm also
\emph{decides} the question, so we keep the test.)
\end{proof}

\subsection{Stage C: constructive minimal-model program}

\begin{lemma}[Finding and contracting $(-1)$-curves]\label{lem:mmp}
There is an algorithm that, given a smooth projective rational surface $S$,
either certifies $S$ is minimal, or finds an explicit $(-1)$-curve $E\subset S$
and computes the contraction $c\colon S\to S'$ (a smooth projective surface) with
its explicit inverse, the blow-up $S=\mathrm{Bl}_{p}S'$ at the computed point
$p=c(E)$.  Iterating, the algorithm terminates in a minimal model $S_{\min}$.
\end{lemma}

\begin{proof}
\emph{Detecting $(-1)$-curves.}  A $(-1)$-curve is an irreducible curve
$E\subset S$ with $p_a(E)=0$ and $E^2=-1$, equivalently (adjunction)
$E\cdot K_S=-1$ and $E^2=-1$.  We enumerate candidate curves of bounded degree
in the fixed projective embedding $S\subset\PP^M$ (the degree of a $(-1)$-curve
in a fixed very ample $H$ is bounded by $H\cdot E$, and the possible classes of
$(-1)$-curves with $H\cdot E\le B$ form a finite computable set because the
N\'eron--Severi lattice $\mathrm{NS}(S)$ is finitely generated and its
intersection form is computable from the resolution history).  Concretely:
\begin{enumerate}
\item Compute a finite generating set of $\mathrm{NS}(S)$ together with the
intersection pairing.  For a rational surface, $\mathrm{NS}(S)\cong
\Pic(S)$ is free of rank $\rho=10-K_S^2$ (since $b_1=0$, $b_2=10-K_S^2$ for a
rational surface), generated by the pullbacks of a hyperplane class of
$\PP^2$/ruling classes and the exceptional classes recorded during the
resolution; the intersection numbers among these are known from the blow-up
combinatorics.  $K_S$ is expressed in this basis via the adjunction formula and
the recorded blow-ups.
\item Solve the finite Diophantine system $E^2=-1$, $E\cdot K_S=-1$, $H\cdot
E>0$, $H\cdot E\le B$ over the lattice for classes $[E]$.  For each solution
class, test effectivity and irreducibility by computing
$h^0(S,\cO_S(E))$ (Riemann--Roch plus the cohomology algorithm of
Lemma~\ref{lem:cohomology}) and inspecting the linear system; a representative
curve is recovered as a vanishing locus of a section.
\end{enumerate}
\emph{Contracting.}  Given $E$ with $E^2=-1$, $p_a(E)=0$, Castelnuovo's
contractibility criterion provides a birational morphism $c\colon S\to S'$ onto a
smooth projective surface contracting exactly $E$ to a point $p$.  Effectively,
$c$ is realised by the linear system $|H'|$ where $H'=mH-(mH\cdot E)E$ chosen so
that $H'\cdot E=0$ and $H'$ remains very ample on $S\setminus E$; the morphism
defined by a basis of $H^0(S,\cO_S(H'))$ is computed, its image $S'$ is the
projective surface defined by the kernel of the corresponding ring map
(elimination), and $c^{-1}$ is the blow-up of $S'$ at the computed point $p$,
again explicit.  Smoothness of $S'$ at $p$ is guaranteed by the criterion and is
verified by the Jacobian test.

\emph{Termination.}  Each contraction drops $\rho(S)=\mathrm{rank}\,
\mathrm{NS}(S)$ by exactly $1$ and increases $K_S^2$ by $1$.  Since $\rho\ge1$
always and a non-minimal surface has $\rho\ge2$, after at most
$\rho(S)-1$ contractions no $(-1)$-curve remains; that surface is minimal,
$S_{\min}$.  For a rational surface $\mathrm{NS}$ is finitely generated with the
stated rank, so the loop terminates in $\le\rho(S)-1$ steps, each step finite by
the bounded enumeration above.
\end{proof}

\begin{remark}[Why minimal $\Rightarrow$ $\PP^2$ or $\mathbb F_e$]
A minimal rational surface $S_{\min}$ (no $(-1)$-curves, $q=P_2=0$) is, by the
Enriques--Castelnuovo classification (Theorem~\ref{thm:mmp}), isomorphic to
$\PP^2$ or to a Hirzebruch surface $\mathbb F_e$, $e\ne1$.  Algorithmically we do
not need to identify $e$ abstractly: we distinguish the two cases by $K_S^2$
($K^2=9$ for $\PP^2$; $K^2=8$ for $\mathbb F_e$) and by the presence of a base-point-free
pencil $|F|$ with $F^2=0$, $F\cdot K=-2$ (a ruling), which exists iff
$S_{\min}\cong\mathbb F_e$ and is found by the same linear-system search as in
Lemma~\ref{lem:mmp}.
\end{remark}

\subsection{Stage D: parametrizing the minimal model}

\begin{lemma}[Parametrizing $\PP^2$ and $\mathbb F_e$]\label{lem:param-min}
There is an algorithm producing an explicit birational map
$\gamma\colon\PP^2_k\dashrightarrow S_{\min}$.
\end{lemma}

\begin{proof}
\emph{Case $S_{\min}=\PP^2$.}  Then $\gamma=\mathrm{id}$ (after expressing the
computed embedding of $S_{\min}\subset\PP^M$ as the $2$-uple--type image of
$\PP^2$; the very ample system used to embed $S_{\min}$ is a multiple of the
hyperplane class of $\PP^2$, and reading the embedding backwards gives explicit
coordinates $[s:t:u]\mapsto S_{\min}$, the Veronese-type map, which is rational
and birational onto its image).

\emph{Case $S_{\min}=\mathbb F_e$.}  Use the ruling $p\colon S_{\min}\to\PP^1$
found above; its generic fibre is a smooth conic over the field
$K=k(\PP^1)=k(\lambda)$ (after restricting to a section's complement one
realises $S_{\min}$ as a conic bundle).  Concretely the bicanonical / ruling
data express $S_{\min}$ birationally as the closure of an affine conic
\[
Q(\lambda;y,z)=a(\lambda)y^2+b(\lambda)yz+c(\lambda)z^2+d(\lambda)y+e(\lambda)z+f(\lambda)=0,
\qquad a,\dots,f\in k[\lambda],
\]
computed by elimination from the equations of $S_{\min}$ relative to $p$.  A
$k(\lambda)$-rational point of this conic gives, by stereographic projection
through that point, a birational parametrization of the conic over $k(\lambda)$,
hence of $S_{\min}$ over $\PP^1$, i.e.\ a birational map
$\AA^1_{\lambda}\times\AA^1_{\mu}\dashrightarrow S_{\min}$, which is the desired
$\gamma$ (with $(s,t)=(\lambda,\mu)$).  A $k(\lambda)$-rational point exists and
is computable: a smooth conic over $k(\lambda)$ with $k$ algebraically closed of
characteristic $0$ has a $k(\lambda)$-rational point by Tsen's theorem (the field
$k(\lambda)$ is a $C_1$ field, so every quadratic form in $\ge3$ variables over
it is isotropic).  The point is found algorithmically by Tsen's effective
procedure: writing the conic as a ternary quadratic form
$\sum q_{ij}(\lambda)w_iw_j=0$ with $q_{ij}\in k[\lambda]$ of degrees $\le D$,
search for a nonzero isotropic vector $w=(w_0,w_1,w_2)$ with $w_i\in k[\lambda]$,
$\deg w_i\le D$; substituting unknown coefficients and equating the resulting
polynomial in $\lambda$ to zero produces a system of polynomial equations over
$k$ which, by Tsen's theorem, has a solution; since $k$ is algebraically closed
the solution is found by elimination (Gr\"obner) and root extraction in $k$.
Stereographic projection from the found point $P_0$ gives explicit rational
functions parametrizing the conic, hence $\gamma$.
\end{proof}

\subsection{Stage E: assembly}

\begin{theorem}[Main algorithm]\label{thm:main}
Algorithm~\ref{alg:main} terminates on every input as in
Definition~\ref{def:size}; it outputs \textsc{Not parametrizable} iff $X$ is not
rational, and otherwise outputs $\varphi\in k(s,t)^{n+1}$ defining a dominant
rational map $\PP^2_k\dashrightarrow X\subset\PP^n_k$ which is birational onto
$X$ (hence a rational parametrization with Zariski-dense image), or, in the
affine input case, $\varphi\in k(s,t)^{n}$ with dense image in
$X^{\mathrm{aff}}$.
\end{theorem}

\begin{algorithm}
\caption{Rational parametrization of a surface}\label{alg:main}
\begin{algorithmic}[1]
\Require Saturated homogeneous ideal $I(X)\subset k[x_0,\dots,x_n]$ of an
irreducible projective surface $X$ (affine inputs replaced by projective closure,
Def.~\ref{def:size}).
\Ensure \textsc{Not parametrizable}, or $\varphi(s,t)\in k(s,t)^{n+1}$ birational
$\PP^2\dashrightarrow X$.
\State Compute a smooth projective model $S$ and birational $\rho\colon S\to X$,
$\rho^{-1}$ \Comment{Lemma~\ref{lem:resolution}}
\State Compute $q(S)=h^1(\cO_S)$ and $P_2(S)=h^0(2K_S)$ \Comment{Lemma~\ref{lem:cohomology}}
\If{$q(S)\ne0$ \textbf{or} $P_2(S)\ne0$}
  \State \Return \textsc{Not parametrizable} \Comment{Prop.~\ref{prop:decision}}
\EndIf
\State $\beta\gets\mathrm{id}$; $T\gets S$
\While{$T$ has a $(-1)$-curve $E$} \Comment{Lemma~\ref{lem:mmp}}
  \State Compute contraction $c\colon T\to T'$ and its inverse blow-up
  \State $\beta\gets c\circ\beta$; $T\gets T'$
\EndWhile
\State $S_{\min}\gets T$ \Comment{minimal; $\PP^2$ or $\mathbb F_e$}
\State Compute birational $\gamma\colon\PP^2\dashrightarrow S_{\min}$
\Comment{Lemma~\ref{lem:param-min}}
\State $\varphi\gets\rho\circ\beta^{-1}\circ\gamma$ \Comment{compose explicit
rational maps; clear denominators}
\State Dehomogenize source: set $[s:t:1]$ to obtain $\varphi(s,t)$
\State \Return $\varphi(s,t)$
\end{algorithmic}
\end{algorithm}

\begin{proof}[Proof of Theorem~\ref{thm:main}]
\emph{Termination.}  Stage A terminates by Lemma~\ref{lem:resolution}
(Zariski/Hironaka resolution, finitely many blow-ups).  Stage B terminates by
Lemma~\ref{lem:cohomology} (finite Gr\"obner computations).  The
\textbf{while} loop terminates by Lemma~\ref{lem:mmp} after $\le\rho(S)-1$
iterations.  Stage D terminates by Lemma~\ref{lem:param-min} (finite
elimination and root finding, justified by Tsen's theorem in the $\mathbb F_e$
case).  Composition and dehomogenization in Stage E are finite polynomial
arithmetic.

\emph{Correctness of the decision.}  By Proposition~\ref{prop:decision}, the
algorithm outputs \textsc{Not parametrizable} exactly when $X$ is not rational.

\emph{Correctness of the parametrization.}  Suppose $q(S)=P_2(S)=0$, so $X$ is
rational and the algorithm proceeds.  Each map in the chain
\[
\PP^2 \xrightarrow{\ \gamma\ } S_{\min}
\xleftarrow{\ \beta\ } S \xrightarrow{\ \rho\ } X
\]
is birational: $\gamma$ by Lemma~\ref{lem:param-min}, $\beta$ as a finite
composition of blow-downs (each birational) by Lemma~\ref{lem:mmp}, and $\rho$
by Lemma~\ref{lem:resolution}.  Therefore the composite
$\varphi=\rho\circ\beta^{-1}\circ\gamma\colon\PP^2\dashrightarrow X$ is
birational, in particular dominant with Zariski-dense image.  Birationality
means $\varphi$ induces an isomorphism $k(X)\xrightarrow{\sim}k(\PP^2)=k(s,t)$,
which both confirms $k(X)\cong k(s,t)$ and gives the explicit parametrization.
Each constituent map was produced as an explicit tuple of homogeneous
polynomials (resp.\ rational functions); their composition is computed by
substitution and denominator clearing, yielding $\varphi$ with entries in
$k(s,t)$.  Restricting the source to the affine chart $[s:t:1]$ gives
$\varphi(s,t)$.  In the affine-input case we further set $x_0=1$ in the target
and discard the $0$-th coordinate, obtaining $\varphi\in k(s,t)^n$ dominating
$X^{\mathrm{aff}}$, because $X^{\mathrm{aff}}$ is dense in its projective closure
$X$ and $\varphi$ has dense image in $X$.
\end{proof}

\section{Complexity analysis}

We measure complexity in the input parameters $n,m,d,H$ (Def.~\ref{def:size}),
counting field operations in the effective subfield of $k$ and bounding the
height growth.

\begin{lemma}[Gr\"obner cost]\label{lem:groebner}
A Gr\"obner basis of an ideal generated by polynomials of degree $\le d$ in $\nu$
variables can be computed in time bounded by $d^{\,2^{O(\nu)}}$ field operations,
and the degrees of the output polynomials are bounded by the
Castelnuovo--Mumford/Bayer--Mumford bound $(2d)^{2^{\nu-1}}$.  Coefficient
heights grow at most singly exponentially in this degree bound.
\end{lemma}

\begin{proof}
These are the standard worst-case bounds for Gr\"obner bases over a field
(Bayer--Mumford; M\"oller--Mora; Dub\'e's degree bound
$2\big(\tfrac{d^2}{2}+d\big)^{2^{\nu-1}}$).  Linear algebra on the graded pieces
up to the degree bound dominates, giving the stated time.
\end{proof}

\begin{theorem}[Overall complexity]\label{thm:complexity}
Algorithm~\ref{alg:main} runs in time bounded by a function of the form
\[
T(n,m,d,H)\ \le\ \big(d\cdot H\cdot m\big)^{\,2^{O(n)}},
\]
i.e.\ doubly exponential in the number of variables and polynomial-in-the-
exponent in $d,H,m$; equivalently the running time is \emph{elementary recursive}
in the input size and singly exponential in the data once the ambient dimension
$n$ is fixed.  The bottleneck is the family of Gr\"obner-basis computations in
Stages A, B, C, D.
\end{theorem}

\begin{proof}
We bound each stage; throughout, all polynomial data stays in $\le n+O(1)$
variables of degree $\le D^*:=(2d)^{2^{O(n)}}$ (Lemma~\ref{lem:groebner}), with
heights bounded by $H^*:=H\cdot D^{*\,O(1)}$.

\emph{Stage A (resolution).}  Each blow-up + normalization is a constant number
of Gr\"obner/elimination/integral-closure computations, each of cost
$D^{*\,2^{O(n)}}$ by Lemma~\ref{lem:groebner}.  The number of blow-ups needed for
surface resolution is bounded by a function of the multiplicity and dimension,
itself bounded by $D^*$ (the resolution invariant strictly decreases and is
bounded by the input degree data).  Hence Stage A costs $D^{*\,2^{O(n)}}=
(2d)^{2^{O(n)}}$.

\emph{Stage B (cohomology).}  Computing a free resolution of the coordinate ring
and the $\mathcal{E}xt$/local-cohomology modules is a bounded number of
Gr\"obner/syzygy computations on modules with the same degree bounds, cost
$D^{*\,2^{O(n)}}$.  Reading off $h^1(\cO_S)$ and $P_2(S)$ is linear algebra over
$k$ on graded pieces of dimension $\le D^{*\,n}$, cost $D^{*\,O(n)}$.

\emph{Stage C (MMP).}  The loop runs $\le\rho(S)-1\le b_2(S)\le D^{*\,O(n)}$
times; each iteration enumerates $(-1)$-curve classes (a lattice problem in rank
$\rho\le D^{*\,O(n)}$, solved by bounded search up to the degree bound) and
performs a contraction (one elimination/Gr\"obner computation). Total cost
$D^{*\,2^{O(n)}}$.

\emph{Stage D (minimal model).}  Elimination to a conic bundle, then Tsen's
isotropic-vector search: the conic has coefficients in $k[\lambda]$ of degree
$\le D^*$, and the resulting polynomial system over $k$ has $O(D^*)$ unknowns and
degree $O(D^*)$, solvable by Gr\"obner + univariate factorization in
$D^{*\,2^{O(1)}}$ field operations; root extraction adjoins algebraic numbers of
degree $\le D^*$.  Cost $D^{*\,2^{O(n)}}$.

\emph{Stage E (assembly).}  Composition of $O(\rho)$ explicit maps of degree
$\le D^*$ multiplies degrees, giving final entries of degree $\le D^{*\,O(\rho)}=
D^{*\,D^{*\,O(n)}}$; this is the only place where degrees can compound, but it is
still bounded by a fixed tower and hence elementary recursive.  (In practice
Stage E degrees are far smaller; the worst-case tower comes from naive
composition and can be reduced by simplifying each map before composing.)

Summing, the dominant terms are the Gr\"obner computations, giving
$T\le(d\,H\,m)^{2^{O(n)}}$ for the decision and reduction stages, with Stage E
contributing the only super-singly-exponential (but still elementary recursive)
factor to the size of the output parametrization.  This proves the bound.
\end{proof}

\begin{remark}[On efficiency]
For fixed ambient dimension $n$ the cost is singly exponential in $d,H,m$, which
is essentially optimal in the worst case because even computing the singular
locus and a single Gr\"obner basis is singly exponential in $d$ for fixed $n$
(Mayr--Meyer-type lower bounds for ideal membership in growing dimension).  For
surfaces of bounded degree (quadrics, cubics, etc.) all bounds become polynomial
in the coefficient height, recovering the classical efficient parametrization
algorithms (e.g.\ stereographic projection for quadrics, the
Schicho/Sendra--Sevilla--Villarino conic-bundle methods) as special cases.
\end{remark}

\section{Verification of the worked invariants}

We record the two numerical sanity checks underlying the algorithm.

\begin{lemma}[Invariants of the terminal cases]\label{lem:invariants}
For the minimal models we have:
$\PP^2$: $q=0,\ P_2=0,\ K^2=9,\ \rho=1$;\quad
$\mathbb F_e$: $q=0,\ P_2=0,\ K^2=8,\ \rho=2$.
In both cases $-K$ is big and the surface is rational, consistent with
Castelnuovo's criterion ($q=P_2=0$).
\end{lemma}

\begin{proof}
For $\PP^2$, $K_{\PP^2}=\cO(-3)$, $K^2=9$, $h^i(\cO)=\delta_{i0}$ so $q=0$, and
$2K=\cO(-6)$ has no global sections so $P_2=0$; $\Pic=\mathbb Z$ so $\rho=1$.
For $\mathbb F_e=\PP(\cO\oplus\cO(e))$ over $\PP^1$, $\Pic=\mathbb Z C_0\oplus
\mathbb Z F$ with $F^2=0$, $C_0^2=-e$, $C_0\cdot F=1$, so $\rho=2$; the canonical
class is $K=-2C_0-(e+2)F$, giving $K^2=8$; $h^1(\cO)=0$ (a $\PP^1$-bundle over
$\PP^1$) so $q=0$; and $2K$ is anti-effective so $P_2=0$.  All values are as
stated; in particular $K^2$ distinguishes the two cases, justifying the
case-split in Stage~D.
\end{proof}

\begin{remark}
The equality $\rho=10-K^2$ for rational surfaces (used in
Lemma~\ref{lem:mmp}) follows from Noether's formula
$12\chi(\cO_S)=K^2+c_2(S)$ with $\chi(\cO_S)=1$ (as $q=p_g=0$) and
$c_2(S)=2-4q+b_2=2+b_2=2+\rho$ (as $b_1=0$, $\mathrm{NS}=H^2$ for rational
surfaces), giving $12=K^2+2+\rho$, i.e.\ $\rho=10-K^2$.  This is the bound that
makes the MMP loop terminate in $\le\rho-1=9-K^2$ steps from a smooth model.
\end{remark}

\end{document}
